\documentclass[UTF8,12pt,a4paper]{ctexart}

\usepackage{amsmath,amssymb}
\usepackage{booktabs}
\usepackage{geometry}
\usepackage{graphicx}
\usepackage[hidelinks]{hyperref}
\geometry{left=2.5cm,right=2.5cm,top=2.5cm,bottom=2.5cm}
\newcommand{\keywords}[1]{\par\noindent\textbf{关键词：}#1}

% 这是无官方 LaTeX 模板时的构建基线，不替代当届官方规则。
\title{预热阶段药材温度与含水率的变化}
\author{}
\date{}

\begin{document}
\maketitle

\begin{abstract}
针对药材预热阶段的径向状态变化，采用一维圆柱导热与非线性水分扩散模型，
以烘房环境数据构造表面边界输入，并用有限体积法求解。
本文件保留第一问的模型、计算设置及已有结果，供后续正文组织与排版使用。
\keywords{药材干燥；径向传递；有限体积法}
\end{abstract}

\clearpage
% 本文件可迁入文字手的整篇项目，章号由主文档决定。
% 后续由 Sol 补足说明与衔接；公式、表格、算法口径沿用核验后交付版本。
\section{问题一：预热阶段药材温度与含水率的变化}
\label{sec:q1}

\subsection{问题分析与模型简化}
% 扩写任务：从空间分布需求、表面对流及内部传递解释模型选择。
% 合并旧稿“问题一的分析”与“模型建立”引言，避免重复背景。
问题一要求计算前 $1800\,\mathrm{s}$ 内不同径向位置的温度和干基含水率。
药材干燥涉及内部水分迁移和表面交换过程\cite{wang2024}。
以半径为特征长度，热 Biot 数为 $h_TR/k\approx1.3889$，
因此保留内部导热与表面对流热阻。
本文借鉴圆柱物料干燥的一维扩散方程和有限体积离散框架\cite{dasilva2014}。

模型描述固定半径圆柱中段的径向响应，忽略端面影响；
将外部水分输入解释为题目给定尺度下的等效浓度；
忽略蒸发潜热对温度场的反馈。
本问热物性为常数，水分扩散系数仅依赖含水率，两个场可分别求解。
% 不将这些简化改写为已被实验验证，也不声称潜热影响很小。

\subsection{径向传热与水分扩散模型}
\subsubsection{变量与物性参数}
设温度和干基含水率分别为 $T(r,t)$ 与 $C(r,t)$，
$0\le r\le R$；$C$ 为水质量与干物质质量之比。
取 $L=0.25\,\mathrm{m}$、$R=0.02\,\mathrm{m}$，以及
\begin{equation}
\rho=820\,\mathrm{kg/m^3},\qquad
c_p=2600\,\mathrm{J/(kg\cdot K)},\qquad
k=0.36\,\mathrm{W/(m\cdot K)}.
\label{eq:q1-thermal-properties}
\end{equation}
\begin{equation}
D(C)=7\times10^{-9}\exp\!\left(-\frac{0.89}{C}\right)
\quad\mathrm{m^2/s}.
\label{eq:q1-diffusivity}
\end{equation}

\subsubsection{控制方程}
根据能量守恒及傅里叶定律，温度满足
\begin{equation}
\rho c_p\frac{\partial T}{\partial t}
=\frac{1}{r}\frac{\partial}{\partial r}
\left(rk\frac{\partial T}{\partial r}\right).
\label{eq:q1-heat}
\end{equation}
干物质骨架静止且单位总体积内干物质量保持不变时，
恒定干物质体积密度 $\rho_d$ 可从水分守恒式两侧约去，得到
\begin{equation}
\frac{\partial C}{\partial t}
=\frac{1}{r}\frac{\partial}{\partial r}
\left[rD(C)\frac{\partial C}{\partial r}\right].
\label{eq:q1-moisture}
\end{equation}
两式适用于 $0<r<R$、$0<t\le1800\,\mathrm{s}$。
式\eqref{eq:q1-moisture}保留扩散通量的散度形式，
$\rho_d$ 不直接等同于热模型中的 $\rho$。
% 可从旧稿补充含 rho_d 的守恒式，解释约去条件；不添加未知密度数值。

\subsubsection{初始条件与边界条件}
记 $T_s=T(R,t)$、$C_s=C(R,t)$，
$T_\infty(t)$ 与 $C_\infty(t)$ 为外部环境输入。
初态及轴线对称条件为
\begin{equation}
T(r,0)=28\,{}^\circ\mathrm C,\qquad C(r,0)=2.55;
\qquad T_r(0,t)=C_r(0,t)=0.
\label{eq:q1-initial-axis}
\end{equation}
以径向向外为正，表面条件为
\begin{align}
-kT_r(R,t)&=h_T[T_s-T_\infty(t)],
&h_T&=25\,\mathrm{W/(m^2\cdot K)},\label{eq:q1-heat-robin}\\
-D(C_s)C_r(R,t)&=h_m[C_s-C_\infty(t)],
&h_m&=8\times10^{-7}\,\mathrm{m/s}.\label{eq:q1-mass-robin}
\end{align}
表面状态由内部传递与外部交换共同决定。
$-D(C)C_r$ 为约化通量，乘以 $\rho_d$ 后才对应水质量通量。
% 补充热边界负号的物理解释；不把 T_s 直接设为 T_infty。

\subsection{环境输入与数值求解}
\subsubsection{环境数据的分段线性插值}
采用附件1中 $0$ 至 $1800\,\mathrm{s}$ 的31个环境节点。
以 $B$ 代表 $T_\infty$ 或 $C_\infty$，相邻采样点间取
\begin{equation}
B(t)=B(\tau_j)+\frac{t-\tau_j}{\tau_{j+1}-\tau_j}
[B(\tau_{j+1})-B(\tau_j)],\quad \tau_j\le t\le\tau_{j+1}.
\label{eq:q1-interpolation}
\end{equation}
% 说明环境采样间隔60 s、积分步长与结果输出间隔1 s是不同概念。

\subsubsection{有限体积离散与非线性迭代}
设 $\Delta r=R/N$、$r_i=(i-1/2)\Delta r$，
圆柱控制体的体积和界面面积为
\begin{equation}
V_i=\pi L(r_{i+1/2}^2-r_{i-1/2}^2),\qquad
A_{i\pm1/2}=2\pi Lr_{i\pm1/2}.
\label{eq:q1-cell-geometry}
\end{equation}
内部界面扩散系数采用调和平均。
定义 $q=-kT_r$、$f=-D(C)C_r$，
令 $m_T=\rho c_p$、$m_C=1$，$F_T=q$、$F_C=f$。
第一步采用后向欧拉，从第二步开始采用 BDF2：
\begin{equation}
m_\phi V_i\frac{3\phi_i^{n+1}-4\phi_i^n+\phi_i^{n-1}}{2\Delta t}
=A_{i-1/2}F_{i-1/2}^{n+1}-A_{i+1/2}F_{i+1/2}^{n+1}.
\label{eq:q1-bdf2}
\end{equation}
水分方程在每步内更新 $D(C)$ 并作 Picard 迭代，线性子问题用追赶法求解。
取 $N=6400$、$\Delta r=3.125\times10^{-6}\,\mathrm{m}$、
$\Delta t=0.0078125\,\mathrm{s}$，迭代容差为 $10^{-10}$。
轴线、内部及表面的输出重构见附录\ref{app:q1-numerics}。
% 从旧稿补充最大迭代差包含全部单元与表面含水率。
% 本段仅叙述已有计算，不再次运行求解。

\subsection{计算结果及规律分析}
\subsubsection{径向温度分布}
表\ref{tab:q1-temperature}列出题目指定时刻和位置的温度。
\input{tables/temperature}
% 接续解释：1800 s时表面36.7856 ℃、中心33.5753 ℃，温差3.2102 ℃；
% 烘房41.513 ℃，表面仍有响应滞后。沿用旧稿未舍入数值所得差值。

\subsubsection{径向含水率分布}
表\ref{tab:q1-moisture}列出对应的干基含水率。
\input{tables/moisture}
% 接续解释：表层失水明显，1800 s表面1.5102 kg/kg，较初值下降约40.78%；
% 中心2.5500是四位小数显示结果，不表示内部没有迁移。

\subsubsection{温度与水分响应的差别}
% 补写一段有数值依据的解释：alpha约1.6886e-7 m²/s，D(2.55)约4.9377e-9 m²/s，
% 二者约34.20倍。解释径向传播尺度差别，不将其称为实际干燥速度比。
两张表合计70个数值。已有逐秒、径向每 $0.1\,\mathrm{cm}$
的完整结果另存于 \texttt{result1.xlsx}。

\subsection{数值可靠性与适用范围}
网格与时间步细化的既有比较结果见表\ref{tab:q1-refinement}。
\begin{table}[htbp]
\centering
\caption{网格与时间步细化的最大绝对差}
\label{tab:q1-refinement}
\small
\begin{tabular}{llll}
\toprule
检验 & 对比设置 & 温度差/$^\circ\mathrm C$ & 含水率差/(kg/kg) \\
\midrule
空间细化 & $N:3200\to6400$ & $8.353721\times10^{-8}$ & $1.174681\times10^{-5}$ \\
时间细化 & $\Delta t:1/64\to1/128\,\mathrm{s}$ & $1.455303\times10^{-8}$ & $3.686571\times10^{-7}$ \\
\bottomrule
\end{tabular}
\end{table}

空间比较固定 $\Delta t=0.0078125\,\mathrm{s}$，时间比较固定 $N=3200$。
比较覆盖每个场37800个完整输出点，两组比较在论文35个取值点上的四位小数均保持一致。
温度结果与800项圆柱导热级数解的最大差为 $2.746\times10^{-8}\,{}^\circ\mathrm C$。
% 可用旧稿简述级数解与全局守恒检验，避免展开全部求解历史。
% 不声称完整输出所有舍入末位获得严格真值保证。

这些检验支持所用连续模型的数值一致性，不构成实验验证。
模型适用于既定假设下的圆柱中段径向响应；端面影响、收缩及潜热反馈未纳入。
% 本问适用范围到此为止，不生成第二问的模型或结果。


\begin{thebibliography}{9}
\bibitem{wang2024}
王乐意，李长河，刘明政，等.
中药材干燥技术与装备研究现状[J]. 农业工程学报，2024，40(2)：1--28.
DOI: \href{https://doi.org/10.11975/j.issn.1002-6819.202306104}{10.11975/j.issn.1002-6819.202306104}.
\bibitem{dasilva2014}
Da Silva W P, e Silva C M D P S, Gama F J A.
Estimation of thermo-physical properties of products with cylindrical shape during drying:
The coupling between mass and heat[J]. Journal of Food Engineering, 2014, 141:65--73.
DOI: \href{https://doi.org/10.1016/j.jfoodeng.2014.05.010}{10.1016/j.jfoodeng.2014.05.010}.
\end{thebibliography}

\clearpage
\appendix
\section{第一问的边界离散与守恒关系}
\label{app:q1-numerics}
% 补足以下公式的衔接说明；技术细节沿用核验后正文第3.1与第4节。
内部界面取
\begin{equation}
D_{i+1/2}=\frac{2D(C_i)D(C_{i+1})}{D(C_i)+D(C_{i+1})}.
\end{equation}
单元未知量按中心点值解释。轴线处面积为零，内侧通量项自然消失；
输出轴线值采用偶二次函数重构：
\begin{equation}
\phi(0,t)\approx\frac{9\phi_1(t)-\phi_2(t)}8.
\end{equation}
内部输出位置用相邻中心点线性插值。设最外层中心 $r_P=R-\Delta r/2$，定义
\begin{equation}
\ell_R=R\ln\frac{R}{r_P}.
\end{equation}
热边界半单元采用
\begin{equation}
q_R=\frac{T_P-T_\infty}{\ell_R/k+1/h_T},\qquad
T_s=T_\infty+q_R/h_T.
\end{equation}
水分边界的扩散系数积分均值用 Simpson 近似：
\begin{align}
\overline D_R&=\frac{D(C_P)+4D((C_P+C_s)/2)+D(C_s)}6,\\
f_R&=\frac{C_P-C_\infty}{\ell_R/\overline D_R+1/h_m},\qquad
C_s=C_\infty+f_R/h_m.
\end{align}
其依据为局部半单元的准稳态通量关系
\begin{equation}
\int_{C_s}^{C_P}D(c)\,\mathrm dc\approx\ell_Rh_m(C_s-C_\infty).
\end{equation}
内部场量仍按瞬态有限体积方程推进。

相邻控制体共享界面通量。BDF2 步的全局温度守恒关系为
\begin{equation}
\rho c_p\sum_iV_i\frac{3T_i^{n+1}-4T_i^n+T_i^{n-1}}{2\Delta t}
=-A_Rq_R^{n+1},\qquad A_R=2\pi RL.
\end{equation}
水分具有相应的约化守恒形式，第一步按后向欧拉计算残差。
已有计算的最大热量绝对平衡残差为 $2.721\times10^{-14}\,\mathrm W$，
约化水分绝对平衡残差为 $2.027\times10^{-22}\,\mathrm{m^3/s}$；
后者乘以 $\rho_d$ 后才是实际水质量平衡残差。

\end{document}
